% Cover letter for the APR: the executive summary on SPHS letterhead, signed by Jim and
% placed on top of the APR packet (this letter, then apr.pdf). A separate document from
% the CV: compile it on its own (Overleaf: Menu > Main document > cover-letter.tex).
% Update the text each review cycle; counts that grow with the record are printed with
% \cvnumber (data/cv-numbers.tex, from scripts/generate.py).
%
% Sources for 2025--2026: Jim's prose draft (chat, 2026-11-24), revised with him so each
% section ends on recognition, and his correction that teaching includes supervision
% (so the course paragraph says "course teaching"). Figures from main.tex ("Significant
% Research Contributions", the Chickadee, HLTH 230, funding and supervision notes),
% data/venue-profile.tex (3 of the period's journal articles in health journals),
% courses.bib (all 2025--2026 offerings in 2026), grants.bib, service.bib (ih2026convener,
% ih2027tpc, cui2025gc, health2025council, health2025fusc, the appeals and SPHS
% committees), activities.bib (ih2026ecw, assessdht2025panel, nfrf2026exploration),
% scholar.bib (siek2027grand), awards.bib (gradsupervision2026nominee, cui2025service,
% chi2025reviews, dis2026reviews, rphs2025edir) and supervision.bib. Veen Wong defends on
% 30 November 2026 (Jim): once she has, say she completed her PhD.
%
% The text of the letter. The standalone document (cover-letter.tex) and the APR version
% of the CV (main.tex, \ifAPR) both read this file; letters/aprletter.tex defines the
% aprletter environment for each.
\begin{aprletter}{%
The Interim Director and the Annual Performance Review Committee\\
School of Public Health Sciences\\
University of Waterloo}{Dear Interim Director and Members of the SPHS APR Committee,}

Please find enclosed my annotated Faculty of Health CV for the 2025--2026 review period, which included a sabbatical from December 2024 to August 2025. I used the sabbatical to lay the groundwork for my lab's machine-learning research on qualitative analysis, which led directly to my students' recent work in this area (ICHEC 2025, CUI 2025, GI 2026). My research advances public health through human-centred computing in three programs: scalable tools for qualitative analysis of online health communities; co-designed games for health behaviour change and gender-norm shift; and technologies that help families care for aging parents at home. Across teaching, research and service, the period produced work that others now use and build on, and that peers have recognized.

\aprheading{research}{Research and Funding}

The work reached both health and computing audiences. Of my \cvnumber{review-journal} journal articles in the period, three appeared in health journals (two in \emph{ACM Transactions on Computing for Healthcare}, one in \emph{Social Science \& Medicine}), alongside peer-reviewed papers at ICHEC, CUI and Graphics Interface. Five of the six papers from my qualitative-analysis program (2022--25) are among the world's 10\% most cited for their field and year, led by Robert Gauthier's CHI 2022 paper in the top 2\% (SciVal). Bold Sky, a game co-designed with social work students to prevent intimate partner violence, was presented to practitioners at the International Summit on Ending Gender-Based Violence, and Veen Wong's presentation of it won the EDI-R Presentation Award at the 2025 Research and Practice in Health Sciences Conference. The National Research Council funded my aging-in-place program a second time (\$321K as Co-PI, 2026--28, after \$201K as PI). New awards also came from CIHR (\$100K, Co-applicant), the Graham Seed Fund (\$35K, PI) and the University, alongside my NSERC Discovery Grant, held continuously since 2015 and more than doubled at renewal. A \$1.1M CIHR Project Grant that I lead as Nominated Principal Applicant is under review.

\aprheading{teaching}{Course Teaching}

My course teaching in the period came in 2026, after my return from sabbatical, and centred on renewal; graduate supervision continued throughout. I redeveloped HLTH 230, and all six of its Student Course Perceptions teaching items rose across its two 2026 offerings. I also rebuilt HLTH 615 around human-centred design with the Centre for Extended Learning. For HLTH 230 I designed and built Chickadee, a programming and assessment platform that now runs as a production campus service. Hosted by IST, it uses single sign-on, sends grades to LEARN (an integration the Dean approved) and passed the University's Policy 46 risk and privacy assessments, with independent security testing. In its Spring 2026 pilot, 143 students made 8,904 submissions and got results in a median of 7.5 seconds. Chickadee is open source, SPHS colleagues plan to adopt it, and a Teaching Innovation Incubator proposal in preparation would bring it to more than 1,300 students in seven Health and Engineering courses in 2027--28. I present it at WatITis, the campus IT conference, on December 9.

\aprheading{teaching}{Graduate Supervision}

I funded nine of my ten graduate students from my grants, and trainees first-authored \cvnumber{review-works-trainee-first} of my \cvnumber{review-works} papers and presentations in the period. Three MMath students completed within program limits and moved into roles including Layer 6 (TD), and Veen Wong defends her PhD on November 30. I sat on 11 graduate committees in 2025 and 12 in 2026. My students nominated me for the University's Award of Excellence in Graduate Supervision.

\aprheading{service}{Service and Leadership}

I led two major international initiatives. As General Chair of ACM CUI 2025, I brought an international conference of about 130 attendees to Waterloo and received its Recognition of Service. As a founding convener of the ACM Conference on Interactive Health, I helped establish a new ACM venue for HCI and health research. Its first conference (Porto, July 2026) drew about 230 attendees in person; I co-chaired its Early Career Workshop for 15 early-career researchers and am its 2027 Technical Program Chair, overseeing peer review of its archival papers. I also co-organized ``Grand Challenges in HCI~\texttimes{} Health,'' a workshop accepted for CHI 2027. In governance, I served as Chair Elect (2025--26) and now Chair (2026--27) of the Health Faculty Council, and on the Faculty Undergraduate Studies Committee, two Faculty appeals committees and three SPHS program committees. I also served on the EU ASSESS-DHT expert panel on digital health technology assessment and an NFRF review panel, and my peer reviews received special recognition at CHI 2025 and DIS 2026.

The CV's notes place this record in the context of HCI's norms, with field-weighted citation evidence. Because the SPHS guidelines provide for expert opinion where assessing quality and impact calls for domain expertise, I ask that the Director obtain it from a University colleague with a regular appointment in HCI.

\aprclosing{Sincerely,}
\end{aprletter}
